\section{Conclusion}
\label{sec:conclusion}

We introduced \eps, a per-token normalized Shannon entropy over
top-$k$ log-probabilities, AST-filtered and max-aggregated into a
per-function score that also preserves per-token attribution. The
function-level score is a recall-maximizing filter: on the evaluated
set, it flags $97\%$ of \textsc{high} API-integration prompts and
catches $100\%$ of true failures at precision $0.272$. The function-%
level claim on its own is not distinguishing. The sequence-level
mean token probability \papg of Spiess et
al.~\cite{spiess2025calibration} reads the same logprobs, costs the
same $1\times$ in API budget, and is competitive with or better
than \eps on function-level detection for three of four OpenAI
models. \papg deserves that credit.

The distinguishing contribution is what the function-level score
gives up to become a single scalar, and what \eps preserves instead.
Of the $\sim 98$ semantic tokens in a typical flagged function,
$10.5\%$ exceed the flag threshold and $1.5\%$ exceed the paused
threshold. A reviewer given \eps's output is not asked to re-read
$98$ tokens; they are asked to evaluate $10$ at the flag threshold
and $1$--$2$ at the paused threshold --- an $89$--$98.5\%$
reduction in within-function review surface. No sequence-level or
inter-generational signal in the literature produces this
localization at the same cost. \papg cannot: it is one scalar per
function. Multi-sample cannot: a $0.52/0.47$ internal split
produces a unanimous output, as GPT-4o on Stripe Scenario A
empirically demonstrates ($\eps = 0.878$, sample diversity $=0.00$
across $N=5$ at $T=0.7$). UQLM detects $5\%$ of our \textsc{high}
prompts on the same benchmark. The intra-generational signal is a
different quantity from the inter-generational one, and it is the
quantity that happens to be actionable for code.

The two-stage architecture is how this becomes a working system.
Stage $1$ is the recall-maximizing filter; stage $2$ is a parallel
LLM reviewer directed to the specific high-\eps tokens rather than
asked to re-read the function. On $201$ entries across three
models, a context-primed reviewer clears $68.7\%$ of flagged entries
with zero late false positives. The token-level attribution is not
decoration; it is the mechanism that makes the reviewer fast enough
to keep up with machine-speed code generation.

The same mechanism --- early-token commitment forcing later tokens,
leaving the genuine uncertainty concentrated in a $1$--$2\%$ slice
of the generation --- plausibly applies in other structured-output
domains (SQL, configuration files, typed program synthesis);
empirical verification is left to future work. Generation is fast;
review is the bottleneck. A filter that misses nothing, plus a
reviewer that is told exactly where to look, moves review off the
developer's critical path.
